\section{IPI Reference Schema}
\label{app:ipi}

Table~\ref{tab:ipi-schema} separates the application objects from their
serialized profiles. Message-mode calls preserve typed or pre-encoded J2735
content. Operation-mode calls add the association and outcome fields required
by a continuing CAV interaction. The runtime stores active sessions and
responses, while each binding serializes the fields required by its path.

\begin{table*}[t]
  \caption{IPI application objects and serialized profiles.}
  \label{tab:ipi-schema}
  \centering
  \footnotesize
  \setlength{\tabcolsep}{3pt}
  \begin{tabular}{@{}p{0.18\textwidth}p{0.34\textwidth}p{0.39\textwidth}@{}}
    \toprule
    Component & Principal fields or encoding & Function \\
    \midrule
    Application API metadata & Message, source, context, send time, transport class, and optional session, correlation, sequence, priority, and expiration fields & Supplies application identity and context; a serialized profile selects the fields required by its exchange \\
    Compact J2735 content profile & One-byte message type, four-byte big-endian payload length, and complete J2735 bytes & Carries typed BSM, PSM, MAP, SPaT, SRM, and SSM content or preserves a pre-encoded J2735 message in $P+5$~B \\
    Correlated CAV operation profile & Session and vehicle identifiers, service class, request/update/completion/rejection role, option flags, expiration, confidence, and length-delimited typed or opaque content & Associates state and outcomes with a perception, planning, control, or map operation \\
    Runtime session record & Vehicle profile, role, requested operations, state, monotonic lease, heartbeat interval, telemetry, and responses & Maintains local state for registration, heartbeat, invocation, response retrieval, and termination \\
    \bottomrule
  \end{tabular}
\end{table*}

The C++17 implementation maps supported J2735 structures into the IPI
application contract and preserves externally encoded J2735 messages byte for
byte. Multi-byte integers in the deterministic operation and session encodings
use network byte order. Length-prefixed fields are checked before allocation,
and decoders reject truncated records, unsupported enumerations, invalid field
bounds, and trailing bytes. The cooperative-operation expiration field uses
tenths of a second since the epoch. Session leases use a monotonic clock, so a
wall-clock adjustment cannot extend a lease. The measurement adapters record
wall-clock fields for trace alignment and use the sender's monotonic clock for
RTT.

The tests cover typed and pre-encoded J2735 content, independently generated
conformance vectors, operation roles, session lifecycle, correlation, malformed
input, PC5 framing, and broker-backed exchange. The runtime validates session
state, request identity, correlation, expiration, and terminal-response
behavior. Applications and deployments define domain-specific transition
policy, security, admission, radio scheduling, and path selection.

\section{Deployment and Measurement Details}
\label{app:setup}

Table~\ref{tab:radio-details} expands the deployment context from
Section~\ref{sec:setup}. The AirSpeed inventory records versions
\texttt{22.0-24-0.0} and \texttt{22.00-53-0.0}. Its two
enabled n48 cells map to two sector carriers. The primary application traffic
used 70/20/10 and Cell~2. The August 15--19 TDD study records application
direction, serving cell, location, reported signal context, and both cells'
administrative state. Seven overlapping DU Cell exports provide five-minute
cell counters for the available experiment windows. Their zone-less timestamps
are interpreted as EDT from traffic alignment. The CBRS export records
Category~B General Authorized Access
operation, one requested and authorized grant, no suspended grant, and a
successful Spectrum Access System heartbeat.

\begin{table*}[t]
  \caption{Detailed radio, gateway, and endpoint configuration.}
  \label{tab:radio-details}
  \centering
  \scriptsize
  \setlength{\tabcolsep}{2.7pt}
  \begin{tabular}{@{}p{0.16\textwidth}p{0.39\textwidth}p{0.37\textwidth}@{}}
    \toprule
    Element & Retained configuration & Measurement record \\
    \midrule
    5G radio & Dedicated AirSpeed 2900; NR n48; TDD; 40~MHz; 30-kHz control and synchronization-signal subcarrier spacing; 20-ms synchronization block period & Clean n48 band without ambient traffic; one uplink and up to two downlink layers configured; downlink 256-quadrature-amplitude modulation enabled as a cell capability \\
    Cells & Cell~1: absolute radio-frequency channel number 637992 and synchronization signal at 3558.720~MHz; Cell~2: 645334 and 3669.600~MHz & Primary traffic and the matched August 17 blocks used Cell~2; the second August 16 40/40/20 TCP repetition used Cell~1; the 70/20/10 Cell~2 lock was operator-reported \\
    Power and antenna & 33-dBm cell transmit power; 34~dBm/MHz effective isotropic radiated power; 17-dBi cross-polarized integrated antenna & Product pattern approximately $65^{\circ}$ azimuth by $8^{\circ}$ elevation; structural centerline 20~ft above roof level \\
    Duplex configuration & Established application profile: 70/20/10 downlink/uplink/dynamic frames; August 15--19 comparison: 40/40/20 and 70/20/10; fixed 10D4G packing & Route, cell state, and serving cell changed during the August 15--16 transition; August 17 and 19 supplied matched workload comparisons; separate 60/20/20 runs are excluded; the attempted 30/60/10 configuration produced no valid experiment \\
    Vehicle gateway & Cisco Meraki MG52-HW campaign gateway with Telit FN990A40 modem and internal antenna; GL.iNet GL-X3000 replacement gateway; 1-Gbit/s full-duplex Ethernet & The August 19 GL-X3000 repeat records gateway RSRP, Reference Signal Received Quality (RSRQ), and SINR while evaluating both TDD profiles at one stationary location \\
    Session and endpoint & SIM provisioned on Data Network Name \texttt{cisco5g}; operator-reported default 5G QoS Identifier 9; on-site MX250 and directly connected d1 host & Host captures observed TOS \texttt{0x0} for default and \texttt{5qi-mapped} experiment labels \\
    Field context & Sole experiment UE; separate three-drive survey with a Samsung 22 handset running G-NetTrack Pro; 1,178 valid band-48 records; RSRP $-121$ to $-88$~dBm; recovered NR SINR $-20$ to 30~dB & Drive-survey medians: $-108$~dBm RSRP and 9~dB SINR; August 17: operator-reported $-100$-dBm RSRP and $-13$-dB RSRQ; August 18: $-105$ to $-115$~dBm; August 19 gateway medians: $-99$ and $-98$~dBm RSRP and 20-dB SINR \\
    PC5 hardware & Mocar OBU, HUALI RSU, 2020 vendor J2735 stack, packet type \texttt{0x1b}, 4,080-B installed-interface limit & GNSS-joined application outcomes; no received-power, channel, resource-pool, modulation, or retransmission export \\
    \bottomrule
  \end{tabular}
\end{table*}

One CAV gateway serves as the sole physical experiment UE. The dedicated radio
and 40-MHz channel operate in a clean n48 band without ambient contention. The load experiment
places one to four streams behind it. The client-demand experiment places
1--100 closed-loop clients behind the same attachment. These workloads
traverse the gateway, Uu path, private core, transport, and d1 responder. The
matrix excludes public-network handover and roaming. The controlled band
isolates experimentally introduced load and application demand from ambient
radio traffic.
Profile transitions can still change route, cell state, serving cell, signal,
and transport state, so those fields remain part of each comparison.

The 5G radio survey comprises three continuous drives through the testbed.
After removing 123 records from a frozen radio-value interval, it retains
1,178 band-48 records with a median RSRP of $-108$~dBm and a median recovered
NR SINR of 9~dB. We use this survey for route-level coverage and retain the
signal records collected with each stationary application campaign for its
workload comparisons.
The PC5 collections include a reference point, a building-obstructed NLOS
point, successive stationary field points, and four driven routes. Exact GNSS
coordinates are transformed to a local frame in the analysis artifacts.

\section{Workload and Experiment Matrix}
\label{app:matrix}

The July 3--4 PC5 stationary sweep configures complete serialized IPI
performance frames at 0.25, 0.5, 1, and 2~KiB, normally for 1,000 requests per
size at 100-ms intervals and a 1,000-ms timeout. The nonzero settings pad both
the request and response to the configured frame length. A separate zero
control produces a 51--55-B request because the identity and timing header
remains present. Each response copies the request's origin, sequence, initial
send time, and configured length. Timeout-dominant collections retain every transmitted
request. The four route collections use 0.25-KiB serialized IPI frames at
200-ms intervals. One stopped route contains
675 attempts; the other routes contain 1,000 attempts each.

The five stationary uplink-heavy 5G payload campaigns were collected on May 13, 14, 15, 21,
and 22, 2026. Most completed conditions contain 1,000 requests at 200-ms
intervals. Every application request uses an IPI cooperative-service frame.
The requested size denotes the application payload before IPI serialization;
sender logs retain the encoded IPI frame size. The dataset manifest contains
22,431 artifacts from
OpenDAIR-V2X, V2X-Seq, TruckV2X, and V2X-Radar. Images span 25.5--666.9~KiB,
point clouds span 15~KiB--3.85~MiB, and maps span 22.3--44.9~MiB. These ranges
motivate the fixed-size sweeps. The 128--512-KiB and 1--2-MiB transfers
represent chunks drawn from these application-payload scales.

We derive payload sizes from 922 saved V2X-Radar validation outputs. The model
assigns a 256-B header and 96~B per predicted object. Payloads range from
12.9 to 24.4~KiB, with median, p95, and p99 sizes of 19.2, 22.3, and
23.4~KiB. Network replay fills an IPI cooperative-service content field to each
representative size. The
58.6-KiB condition is a separate stress payload. The responder returns a
compact correlated acknowledgment immediately after decoding the frame.

The August 15--16 TDD runs retain the same uplink-heavy timing contract. The
August 17 matched experiment instead uses a compact vehicle request followed by
an exact 1-, 10-, 100-, or 500-KiB edge-to-vehicle response. Each profile
contains 1,000 validated TCP and MQTT exchanges per payload. The vehicle timer
includes complete payload receipt and sequence, length, and checksum
validation. The same experiment then performs ten exact 50-MiB TCP
uploads and ten exact downloads per profile. Endpoint byte counts, checksums,
and receiver-observed durations define application goodput.

The August 18 40/40/20 repeat uses 1,000 uplink-heavy and downlink-heavy
exchanges at 1, 10, and 100~KiB. Each transport contains 250 declared
uplink-heavy exchanges at 500~KiB and 500 downlink-heavy exchanges. Reported RSRP
changes from $-105$ to $-110$ and then $-115$~dBm during the run. Its shortened
bulk stage contains one complete upload/download pair and one additional
upload. The failed August 17 downlink preflight, the bulk pilot, and the
unintended August 18 partial preflight remain diagnostic artifacts outside the
primary estimates.
The August 19 GL-X3000 repeat uses 500 validated exchanges for every
combination of direction, transport, profile, and 1-, 10-, 100-, or 500-KiB
payload. It adds three exact 50-MiB uploads and three exact downloads per
profile. The two profile blocks retain separate gateway radio records and use
the same stationary location.
Table~\ref{tab:experiment-matrix} records the controlled conditions and the
corresponding measurement for every experiment family.

\begin{table*}[t]
  \caption{Complete workload and communication experiment matrix.}
  \label{tab:experiment-matrix}
  \centering
  \scriptsize
  \setlength{\tabcolsep}{2.7pt}
  \begin{tabular}{@{}p{0.18\textwidth}p{0.36\textwidth}p{0.36\textwidth}@{}}
    \toprule
    Question & Controlled conditions & Recorded outcome \\
    \midrule
    Message and compact exchange & J2735 callback checks; local IPI tests; decoded PC5 IPI performance frames with sequence-correlated responses; Uu 0--4-KiB IPI frames over TCP/MQTT & Decode or callback result, response completion, RTT \\
    Payload scale & PC5 minimum-header control and 0.25--2-KiB complete IPI frames; uplink-heavy Uu IPI frames with 0--2~MiB vehicle-originated content; downlink-heavy Uu frames with 1--500~KiB returned content; detector-derived sizes; 58.6-KiB stress payload & $C(D)$ and response p50/p95/p99 \\
    Field and route coverage & PC5 stationary points and four routes; original Uu stationary collections and Samsung 22 continuous-drive survey; matched August 17 location~3; August 18 time-segmented context; August 19 same-location gateway measurements & GNSS-joined response completion, RTT, and separate radio context \\
    Protocol formation & TCP, MQTT QoS~0 over TCP, raw UDP, and UDP with 1,400-B adapter-level fragments & Returned acknowledgment after decode or returned payload after reassembly and validation \\
    TDD and application direction & August 15--16 uplink-heavy repetitions; August 17 matched downlink-heavy responses and exact transfers; August 18 weak-signal direction repeat; August 19 same-device 1--500-KiB bidirectional and exact-transfer matrix under both profiles & Profile-transition repeatability, response RTT, deadline completion, exact application goodput, device sensitivity, and radio context \\
    Exact directional transfer & Ten exact 50-MiB uploads and downloads per August 17 profile; three per August 19 profile and direction; receiver byte-count and checksum validation & Individual transfer goodput, mean, and 95\% confidence interval \\
    Competing traffic & Separate historical 15-minute offered controls; 1-KiB foreground with one to four offered 25-Mbit/s uplink streams & Achieved background rate and foreground deadline completion \\
    Application demand & 1, 2, 5, 10, 20, 50, and 100 clients; 1-KiB closed-loop requests; sole experiment UE & Aggregate and per-client response completion and RTT \\
    Packet marking & Default and \texttt{5qi-mapped} application labels & Vehicle-interface Internet Protocol header capture \\
    Interruption & TCP receiver, UDP receiver, MQTT receiver, or MQTT broker restarted for 10~s & Missing attempts, deadline completion, and response gap \\
    \bottomrule
  \end{tabular}
\end{table*}

In the closed-loop application-demand collection, each application client waits
for a response or timeout before the next
request from that client and then applies a 200-ms interval. At negligible
response time, five requests/s/client is the maximum nominal rate. A
locked-cell 70/20/10 collection additionally uses four conditions: one idle 1-KiB client,
one idle 23.4-KiB client, the 1-KiB client with offered uplink traffic, and 100
1-KiB clients. TCP, MQTT, and UDP occupy separate five-minute collection bins
with two repetitions. The first three conditions use 500 probes per repetition;
the 100-client condition uses 500 probes per client per repetition.

Every issued primary request/response attempt remains in its condition
denominator, including missing, timed-out, late, and failed attempts. Entirely
excluded diagnostic roots and never-started bulk repetitions do not enter a
primary estimate. Operator-shortened conditions use their declared retained
counts. The request/response experiments carry application-derived workloads
in IPI cooperative-service frames and measure RTT and deadline completion.
PC5 uses decoded IPI performance requests and responses matched by origin and
sequence identifiers. Exact bulk transfers instead use
validated bytes, checksum, receiver duration, and application goodput.

\section{Detailed Results for Additional Conditions}
\label{app:detailed-results}

Several experiment families include conditions that are not represented by
dedicated figure panels. This section reports their complete numerical
outcomes.

\subsection{Direct PC5 Route Results}

Each route used a complete 0.25-KiB serialized IPI request. Routes 1 and
2 used a 1,000-ms timeout, while Routes 3 and 4 used a 500-ms timeout. The
response percentiles include only sequence-matched responses; completion uses
all issued requests as its denominator.

\begin{table*}[t]
  \caption{Direct PC5 route response completion and RTT.}
  \label{tab:pc5-route-results}
  \centering
  \scriptsize
  \setlength{\tabcolsep}{4.0pt}
  \begin{tabular}{@{}lrrrrrrr@{}}
    \toprule
    Route & Attempts & Responses & Completion (\%) & Timeout (ms) & p50 RTT (ms) & p95 RTT (ms) & p99 RTT (ms) \\
    \midrule
    R1 & 675 & 515 & 76.3 & 1,000 & 99.610 & 108.421 & 112.436 \\
    R2 & 1,000 & 731 & 73.1 & 1,000 & 28.326 & 41.887 & 55.880 \\
    R3 & 1,000 & 591 & 59.1 & 500 & 29.204 & 43.923 & 54.317 \\
    R4 & 1,000 & 716 & 71.6 & 500 & 28.276 & 41.621 & 50.651 \\
    \bottomrule
  \end{tabular}
\end{table*}

Across the four routes, 59.1--76.3\% of issued requests received a response.
The p95 RTT among received responses was below 50~ms on three routes even
though 26.9--40.9\% of all requests on those routes timed out. The GNSS join
places missing responses in recurring no-response segments.

\subsection{Field-Specific Uplink Payload and Transport Results}

Field collections changed the response tail across both payload and transport.
In one collection, 0--4-KiB p95 RTT was 36--49~ms for MQTT and 100--140~ms for
TCP. At 2~MiB, p95 reached 2,536 and 3,633~ms. At another location, 2-MiB TCP
p95 reached 10,767~ms. In a variable-path collection, empty-payload TCP and
MQTT p95 rose to 1,142 and 510~ms. The 256-KiB MQTT condition returned
999/1,000 responses with a 43,836-ms p95.

The transport comparison used empty-payload controls, detector-derived
payloads, and the 58.6-KiB stress payload. In the $-106$-dBm collection, TCP and
MQTT returned every response at each size. MQTT p95 RTT increased from
49.826~ms with an empty payload to 183.705~ms for the 23--24-KiB group and
403.790~ms at 58.6~KiB. The corresponding TCP values were 151.696,
169.540--213.894, and 381.409~ms. In the $-120$-dBm collection, 58.6-KiB MQTT
and TCP p95/p99 increased to 770/1,230 and 662/1,027~ms.

Raw UDP returned 1,000/1,000 empty-payload and 300/300 1-KiB responses. It
returned 0/45, 0/100, and 0/100 at 1.37, 4, and 19.2~KiB. Application
fragmentation and receiver reassembly increased response completion at
$-106$~dBm. Completed responses reached 99.5--100.0\% through 23.4~KiB, while
the 58.6-KiB condition returned 20/100 responses. At $-120$~dBm, payloads
through 4~KiB returned 99.9--100.0\%, 19--25-KiB payloads returned
93.5--96.1\%, and 58.6~KiB returned 0/100.

\subsection{Gateway and TDD Replication Results}

At 70/20/10, the closest MG52 control and the GL-X3000 repeat had reported or
measured median RSRP of $-100$ and $-99$~dBm. The GL-X3000 repeat measured
15.691~Mbit/s upload and 166.106~Mbit/s download goodput. These rates are
30.8\% and 31.0\% above the MG52 control rates of 12.000 and
126.786~Mbit/s. Under 40/40/20, the earlier MG52 collection crossed weaker
reported signal levels from $-105$ to $-115$~dBm. The later GL-X3000
collection at median $-98$~dBm measured 10--500-KiB vehicle-to-edge p50 RTTs
63.7--92.4\% below the earlier MG52 collection, so the earlier large-uplink
collapse did not recur. The same-device comparison still produced lower burst
latency and higher sustained goodput under 70/20/10. The higher rates in the
replacement-device run therefore did not remove the measured gap between the
two deployed configurations.

Consecutive 70/20/10 runs changed TCP p95 RTT by only 0.22--1.29\% across
1--100~KiB. During the initial 40/40/20 rollout, the route required recovery,
cell administrative state changed, and separate repetitions selected different
serving cells. A later matched MG52 downlink block completed all 8,000
exchanges, although its payload and transport effects differed. The GL-X3000
repeat completed both directional matrices but retained the large-object
latency and goodput gap.

\subsection{Sustained-Stream Load and Packet-Marking Results}

The two $-106$-dBm repetitions achieved 1.320--2.267~Mbit/s of background
traffic. MQTT $C(100\,\mathrm{ms})$ fell from 99.70\% at idle to 40.15\% with
the generator active. At 200~ms, it changed from 99.95\% to 98.00\%. TCP
achieved 0.75\% and 0.40\% at 100~ms. Its 200-ms completion fell from 97.90\%
to 59.00\%. Across one, two, and four background streams, achieved throughput
ranged from 1.990 to 5.743~Mbit/s. The pooled active-generator conditions in
two weak-band collections produced MQTT/TCP 200-ms completion of
83.63\%/38.40\% and 91.90\%/49.17\%.

At $-106$~dBm, the active generator increased MQTT p95 from 52--60 to
174--176~ms. TCP p95 increased from 162--194 to 388--420~ms. Across the one-,
two-, and four-stream sweep, MQTT p95 was 170, 468, and 182~ms. TCP p95 was
391, 402, and 410~ms. The MQTT tail did not vary monotonically with the number
of streams.

The default and \texttt{5qi-mapped} labels both produced observed TOS values of
\texttt{0x0} in the vehicle capture. The foreground and background flows
entered the measured path without an observed packet-class difference between
the requested labels.

\subsection{Concurrent-Client Results}

In the $-106$-dBm seven-point repeat, 100 clients produced aggregate p95 RTTs
of 81.850~ms for MQTT, 264.070~ms for TCP, and 113.391~ms for UDP. MQTT and TCP
each returned 100,000 of 100,000 responses. UDP returned 99,965 of 100,000. A
separate 100-client repeat yielded p95 values of 177.053, 283.765, and
145.894~ms and returned 100,000, 100,000, and 99,365 responses.

The 100-client collection at $-120$~dBm produced a different ordering. MQTT
and TCP p95 rose to 5,135.984 and 2,908.420~ms. They completed 99.499\% and
100\% of issued requests before timeout. UDP completed 90,811 of
100,000 requests, with a p95 of 172.500~ms. Its shorter response tail coincides
with 9,189 missing cycles.

The per-client distributions reveal substantial unfairness within the shared
attachment. Across 100 clients, the median client p95/p99 was
3,814.102/13,149.421~ms for MQTT and 2,938.408/6,400.146~ms for TCP. Many flows
fall outside the 100- and 200-ms application budgets even though most requests
still receive a response before timeout.
